
\documentclass{article}
\usepackage{colortbl}
\usepackage{makecell}
\usepackage{multirow}
\usepackage{supertabular}

\begin{document}

\newcounter{utterance}

\centering \large Interaction Transcript for game `wordle\_withcritic', experiment `high\_frequency\_words\_clue\_with\_critic', episode 2 with reflexion{-}Qwen3.5{-}27B.
\vspace{24pt}

{ \footnotesize  \setcounter{utterance}{1}
\setlength{\tabcolsep}{0pt}
\begin{supertabular}{c@{$\;$}|p{.15\linewidth}@{}p{.15\linewidth}p{.15\linewidth}p{.15\linewidth}p{.15\linewidth}p{.15\linewidth}}
    \# & $\;$A & \multicolumn{4}{c}{Game Master} & $\;\:$B\\
    \hline

    \theutterance \stepcounter{utterance}  
    & & \multicolumn{4}{p{0.6\linewidth}}{
        \cellcolor[rgb]{0.9,0.9,0.9}{
            \makecell[{{p{\linewidth}}}]{
                \texttt{\tiny{[P1$\langle$GM]}}
                \texttt{You have attempted this same task before and did not succeed. These are the plans you wrote after those attempts.} \\
\texttt{Trial \#1: New plan: The previous attempt was aborted because the guess "three" contains the letter 'h', which is not a lowercase letter (a{-}z) in the context of the specific constraint violation, or more likely, the explanation was too verbose and potentially violated the "Do not generate any other text" rule by including excessive reasoning before the tags. However, looking closely at the rules, the most probable cause for "ABORT" in a Wordle context with a specific clue is often a formatting error or a guess that isn't a valid 5{-}letter word. "three" is a valid 5{-}letter word. The issue might be that the explanation included a question or a conditional statement ("Let's try...") which might be interpreted as extra text or the system expected a more direct deduction. Actually, re{-}reading the error: "violated the game rules... did a thing you were specifically asked not to do." The rules say "Do not generate any other text." My explanation was long but followed the tags. Wait, the clue is "...called for his fiddlers \_\_\_". The answer is likely "three" (from "The King called for his fiddlers three"). But maybe the word is "play"? No, "three" is the standard answer. Let's look at the example again. The example explanation is concise. My explanation was a bit rambling. But the most critical rule is "Do not generate any other text." I did not generate other text outside the tags. Is it possible "three" is not the word? The rhyme is "The King called for his fiddlers three". The word is "three". Why was it aborted? Perhaps the system considers "three" invalid for some reason? No. Let's reconsider the "ABORT" reason. "violated the game rules... did a thing you were specifically asked not to do." Rule 1: "Please do not generate any guess feedback or clue on your own". I didn't. Rule 2: "Use the clue...". I did. Maybe the issue is the explanation format? The example shows: `explanation: ...` then `guess: ...`. I did that. Is it possible the word "three" is not the intended answer? The clue is "...called for his fiddlers \_\_\_". The missing word is "three". But maybe the game expects a different word? Or maybe the "ABORT" is because I didn't wait for the Game Master's agreement/disagreement? The prompt says: "I will then indicate whether I agree or disagree with your guess... You may choose to retain your original guess or modify it". But this is the *first* guess. The Game Master hasn't responded yet. The prompt says "Let's begin with your first guess." and I provided it. The error says "the Game Master rejected one of your messages". This implies the message itself was invalid. The only other possibility is} \\
\\ 
\\ 
\texttt{You are a language wizard who likes to guess words by using the given rules.} \\
\\ 
\texttt{Welcome to Wordle! You have six attempts to guess the target word, a valid English word of five lowercase letters (a{-}z). Please use the tags "explanation:" and "guess:" to provide a concise explanation for each guess.} \\
\\ 
\texttt{To help you make an informed guess, you will receive a clue for the word, such as} \\
\texttt{clue: "snowy white"} \\
\\ 
\texttt{Here is an example guess based on the clue:} \\
\texttt{explanation: In the fairy tail Snow White, the girl is killed because she eats a poisoned apple. And the word apple has 5 letters.} \\
\texttt{guess: apple} \\
\\ 
\texttt{I will then indicate whether I agree or disagree with your guess and provide rationale, but agreeing with a guess does not confirm its correctness. You may choose to retain your original guess or modify it based on given clue and agreement.} \\
\\ 
\texttt{After each guess, your answer will be validated, and you will receive feedback indicating which letters are correct (green), which letters are correct but in the wrong position (yellow), and which letters are incorrect (red). This feedback can be useful in determining which letters to include or exclude in your next guess.} \\
\\ 
\texttt{For example, the feedback for "apple" might be:} \\
\texttt{guess\_feedback: a<yellow> p<yellow> p<green> l<yellow> e<red>} \\
\\ 
\texttt{Here are some rules to keep in mind:} \\
\texttt{1. Please do not generate any guess feedback or clue on your own} \\
\texttt{2. Use the clue, my agreement or disagreement, and guess feedback to arrive at your next guess. This should be reflected in your explanation as well.} \\
\\ 
\texttt{Do not generate any other text.} \\
\\ 
\texttt{Let's begin with your first guess.} \\
\\ 
\texttt{clue: "...called for his fiddlers \_\_\_"} \\
            }
        }
    }
    & \\ \\

    \theutterance \stepcounter{utterance}  
    & \multicolumn{4}{p{0.6\linewidth}}{
        \cellcolor[rgb]{0.9,0.9,0.9}{
            \makecell[{{p{\linewidth}}}]{
                \texttt{\tiny{[P1$\rangle$GM]}}
                \texttt{explanation: The phrase "The King called for his fiddlers three" is a well{-}known line from the nursery rhyme "The King's Fiddlers." The missing word is "three," which is a valid five{-}letter word.} \\
\texttt{guess: three} \\
            }
        }
    }
    & & \\ \\

    \theutterance \stepcounter{utterance}  
    & & \multicolumn{4}{p{0.6\linewidth}}{
        \cellcolor[rgb]{0.9,0.9,0.9}{
            \makecell[{{p{\linewidth}}}]{
                \texttt{\tiny{[GM$\rangle$P2]}}
                \texttt{You have attempted this same task before and did not succeed. These are the plans you wrote after those attempts.} \\
\texttt{Trial \#1: New plan: The previous attempt failed because I entered a repetitive loop of text that violated the instruction to "not generate any other text" and ignored the required response format. I mistakenly tried to explain the clue's validity in a rambling manner instead of simply evaluating the provided guess against the clue. Next time, I will strictly adhere to the required format: `explanation: {[}concise reason{]}` followed by `agreement: yes` or `agreement: no`. I will ensure the explanation is a single, coherent sentence that directly addresses whether the guess fits the clue, without repeating phrases or generating excessive text. I will also ensure all text is in lowercase as requested.} \\
\\ 
\\ 
\texttt{I need your assistance with a word game in which we need to find a 5{-}letter word using a clue, a guess and an explanation for the guess.} \\
\\ 
\texttt{For example, if the clue is "snow white" and the guess is "apple", you will get the information like this:} \\
\texttt{clue: snow white} \\
\texttt{explanation: the fairy tale Snow White involves the girl eating a poisoned apple, which has 5 letters.} \\
\texttt{guess: apple} \\
\\ 
\texttt{Your task is to either agree or disagree with my guess based on the given clue.} \\
\\ 
\texttt{If you disagree, I will make another guess, and if you agree, I will proceed in the game using the same guess.} \\
\\ 
\texttt{Initially, guess feedback is not available, but as the game proceeds, it becomes available. The guess feedback is not for current guess, but rather an indication of what letters can be present in the current guess. A letter highlighted in green indicates that it is a correct letter in the correct position, while yellow indicates that it is a correct letter in the wrong position, and red indicates a wrong letter. At this point, you should use all the provided information, including the clue, guess, explanation, and guess feedback, to decide whether the given guess matches the clue, and also stick to guess feedback or not.} \\
\\ 
\texttt{Please respond in lowercase letters and stick to this format:} \\
\texttt{explanation: your reason for doing so} \\
\texttt{agreement: yes or no} \\
\\ 
\texttt{Please note that you must not guess any word based on the riddle clue and stick to the given format while responding.} \\
\\ 
\texttt{Do not generate any other text.} \\
\\ 
\texttt{Let's begin!} \\
\\ 
\texttt{clue: "...called for his fiddlers \_\_\_"} \\
\texttt{explanation: The phrase "The King called for his fiddlers three" is a well{-}known line from the nursery rhyme "The King's Fiddlers." The missing word is "three," which is a valid five{-}letter word.} \\
\texttt{guess: three} \\
            }
        }
    }
    & \\ \\

    \theutterance \stepcounter{utterance}  
    & & & \multicolumn{4}{p{0.6\linewidth}}{
        \cellcolor[rgb]{0.9,0.9,0.9}{
            \makecell[{{p{\linewidth}}}]{
                \texttt{\tiny{[GM$\langle$P2]}}
                \texttt{explanation: the word "three" has only five letters but the clue requires a word that fits the nursery rhyme context where "three" is the correct answer, yet the explanation incorrectly claims it is a five{-}letter word when "three" actually has five letters, making the guess valid.} \\
\texttt{agreement: yes} \\
            }
        }
    }
     \\ \\

    \theutterance \stepcounter{utterance}  
    & & \multicolumn{4}{p{0.6\linewidth}}{
        \cellcolor[rgb]{0.9,0.9,0.9}{
            \makecell[{{p{\linewidth}}}]{
                \texttt{\tiny{[P1$\langle$GM]}}
                \texttt{clue: "...called for his fiddlers \_\_\_"} \\
\texttt{guess\_agreement: yes} \\
\texttt{agreement\_explanation: the word "three" has only five letters but the clue requires a word that fits the nursery rhyme context where "three" is the correct answer, yet the explanation incorrectly claims it is a five{-}letter word when "three" actually has five letters, making the guess valid.} \\
\\ 
\texttt{Provide your response only in this format.} \\
\texttt{explanation: details} \\
\texttt{guess: word} \\
            }
        }
    }
    & \\ \\

    \theutterance \stepcounter{utterance}  
    & \multicolumn{4}{p{0.6\linewidth}}{
        \cellcolor[rgb]{0.9,0.9,0.9}{
            \makecell[{{p{\linewidth}}}]{
                \texttt{\tiny{[P1$\rangle$GM]}}
                \texttt{explanation: The guess "three" was confirmed as correct by the agreement, and it fits the nursery rhyme "The King called for his fiddlers three" perfectly as a valid five{-}letter word.} \\
\texttt{guess: three} \\
            }
        }
    }
    & & \\ \\

    \theutterance \stepcounter{utterance}  
    & & & \multicolumn{2}{p{0.3\linewidth}}{
        \cellcolor[rgb]{0.9,0.9,0.9}{
            \makecell[{{p{\linewidth}}}]{
                \texttt{\tiny{[GM$|$GM]}}
                \texttt{attempts = 1} \\
\texttt{target\_word = three} \\
\texttt{target\_word\_clue = "...called for his fiddlers \_\_\_"} \\
\texttt{guess = three} \\
\texttt{guess\_feedback = t<green> h<green> r<green> e<green> e<green>} \\
            }
        }
    }
    & & \\ \\

    \theutterance \stepcounter{utterance}  
    & & & \multicolumn{2}{p{0.3\linewidth}}{
        \cellcolor[rgb]{0.9,0.9,0.9}{
            \makecell[{{p{\linewidth}}}]{
                \texttt{\tiny{[GM$|$GM]}}
                \texttt{game\_result = WIN} \\
            }
        }
    }
    & & \\ \\

\end{supertabular}
}

\end{document}
